\section{Measurement records as evidential anchors and the structure of theory-level confirmation}
\label{sec:confirmation}

Section~\ref{sec:semanticcore} isolated the quotient structure of admissible reference assignments and
identified fixation commutation as the bridge invariant. The present section uses that bridge to turn the
quotient into a measurement theorem about evidence. Measurement records are evidential anchors; once
that role is explicit, divergence among admissible quotient classes is no longer a semantic curiosity but
a question about whether the theory fixes its own evidential output.

\begin{claimbox}
\textbf{Central hinge.} The quotient theorem identifies which divergences among admissible reference
assignments are standing-bearing, and Section~\ref{sec:semanticcore} proves that fixation commutation
is the bridge invariant for theory-level evidential use. The present section therefore asks a narrower
question: when quotient-relevant divergence survives into comparative evidential roles, does the theory
fix its own evidential output?
\end{claimbox}

\subsection{Measurement records as evidential anchors}

\begin{definition}[Evidential anchor]
A measurement record $r\in\Rset$ is an \emph{evidential anchor} for a diachronic measurement claim
$d\in\Dset$ just in case the theory-level use of $d$ depends on taking $r$ as the referent of $d$ and
then using the induced evidential proposition $\Prop_\rho(d)$ to compare rival hypotheses in
$\Hset$.
\end{definition}

\begin{definition}[Theory-level confirmation]
\label{def:theory-level-confirmation}
A \emph{theory-level confirmation verdict} is a comparative confirmational verdict attributable to the
theory on the fixed scope rather than to an unresolved choice among admissible reference
assignments. Equivalently, a theory-level confirmation verdict is well-defined only if the theory's
admissible use does not leave the comparative evidential role of its measurement records unsettled
across admissible standing-equivalence classes.
\end{definition}

\begin{definition}[Comparative hypothesis ordering induced by a record]
Let $d\in\Dset$ be an evidential claim and let $\rho\in\ARef$ be an admissible reference assignment.
The \emph{comparative hypothesis ordering induced by $d$ under $\rho$} is the function
\[
\Ord_\rho^d(H_i,H_j):=\CR(\Prop_\rho(d))(H_i,H_j)
\qquad (H_i,H_j\in\Hset).
\]
A theory-level confirmation claim built from $d$ is theory-relative only if this ordering is invariant
on the admissible quotient classes relevant to $d$.
\end{definition}

\begin{proposition}[Evidential use of records is assignment-sensitive]
\label{prop:assignment-sensitive}
Let $d\in\Dset$ be a diachronic measurement claim used as an evidential anchor. Then every
admissible reference assignment $\rho\in\ARef$ determines a corresponding comparative hypothesis
ordering $\Ord_\rho^d$ on $\Hset$. Therefore any unresolved variation among admissible quotient
classes is already variation in evidential use unless all such orderings coincide.
\end{proposition}

\begin{proof}
By definition of admissible reference assignment, $\rho(d)$ is a physically realized and admissible
branch-record pair for $d$. By definition of the induced evidential proposition, $\rho(d)$ determines
$\Prop_\rho(d)\in\Eset$. By definition of confirmational role, $\CR(\Prop_\rho(d))$ is a function on
ordered pairs of hypotheses. So $\Ord_\rho^d$ is well-defined for every admissible assignment. If
different admissible quotient classes induce different orderings, then the record is doing different
comparative evidential work under those classes. Hence unresolved variation among admissible
quotient classes is already variation in evidential use unless all orderings coincide.
\end{proof}

\subsection{The semantic hinge}

\begin{lemma}[Two-case exhaustion for evidential reference]
\label{thm:two-case-exhaustion}
\label{lem:dichotomy-admissible-assignments}
Let $d\in\Dset$ be used evidentially. Then exactly one of the following holds:
\begin{enumerate}
    \item all admissible reference assignments lie in a single $\MReq$-class on the evidential use of
    $d$; or
    \item there exist admissible reference assignments whose evidential values for $d$ lie in distinct
    $\MReq$-classes.
\end{enumerate}
No third same-scope case exists.
\end{lemma}

\begin{proof}
Either the admissible assignments agree on the standing-bearing evidential value of $d$, or they do
not. If they do, then they lie in one $\MReq$-class on that use. If they do not, then there exist at
least two distinct $\MReq$-classes. By
Definition~\ref{def:fixed-downstream-substrate} and Proposition~\ref{prop:substrate-consequences},
no third same-scope standing-bearing locus remains.
\end{proof}

\begin{theorem}[No auxiliary parameterization of theory-level confirmation]
\label{thm:no-auxiliary-parameterization}
Let $T$ be a theory and let $\ARef$ be the class of its admissible reference assignments for the fixed
measurement setting. If comparative confirmational verdicts for $T$ vary across admissible
$\MReq$-classes, then those verdicts are not attributable to $T$ as such. They are attributable only to
the pair consisting of $T$ together with the selected quotient class. Consequently, such verdicts do not
constitute theory-level confirmation in the sense relevant to the target explanatory practice.
\end{theorem}

\begin{proof}
By Definition~\ref{def:theory-level-confirmation}, a theory-level confirmation verdict must be
attributable to the theory on the fixed scope rather than to a further unresolved choice among
admissible reference assignments. Suppose instead that distinct admissible $\MReq$-classes yield
different comparative orderings for some evidential claim $d$. Then the comparative support relation
between hypotheses is not fixed by $T$ alone. It varies with the chosen admissible quotient class. In
that case the verdict is attributable only to $(T,[\rho])$ for some chosen class $[\rho]\in\ARef/\MReq$,
not to $T$ as such. Such verdicts are therefore assignment-parameterized rather than theory-level.
\end{proof}

\begin{theorem}[Theory-level confirmation criterion]
\label{thm:theory-level-confirmation-criterion}
For an evidential claim $d\in\Dset$, theory-level confirmation is well-defined if and only if the map
\[
[\rho] \longmapsto \Ord_\rho^d
\]
from $\ARef/\MReq$ to comparative hypothesis orderings is constant. Equivalently, if there exist
admissible assignments whose comparative confirmational roles differ, then no theory-level
confirmation verdict is fixed by the theory as such.
\end{theorem}

\begin{proof}
If the displayed map is constant, then all admissible assignments induce the same comparative
ordering for $d$, so the verdict is attributable to the theory on the fixed scope rather than to a
selected admissible quotient class. Conversely, if the map is not constant, then there exist admissible
assignments inducing different comparative orderings. By Corollary~\ref{cor:confirmational-divergence-quotient-relevant},
that divergence is quotient-relevant and therefore standing-bearing rather than semantic skin. By
Theorem~\ref{thm:no-auxiliary-parameterization}, the resulting verdict is attributable only to the pair
consisting of the theory together with a chosen admissible quotient class. Hence no theory-level
confirmation verdict is fixed by the theory as such.
\end{proof}

\begin{corollary}[Confirmational consequence]
\label{thm:confirmational-consequence}
\label{cor:confirmation-collapse}
If there exist admissible reference assignments that differ in comparative confirmational role, then
confirmation is assignment-relative rather than theory-relative on the fixed scope.
\end{corollary}

\begin{proof}
Immediate from Theorem~\ref{thm:theory-level-confirmation-criterion}.
\end{proof}

\subsection{Fixation as the bridge invariant}

Definitions~\ref{def:lawful-redescription-evidential} and \ref{def:fixation-commutation}, together with
Theorem~\ref{thm:fixation-commutation-necessity} and
Corollary~\ref{cor:fixation-central-hinge}, were established already in
Section~\ref{sec:semanticcore}. They are placed there because fixation is the bridge invariant between
quotient structure and all later applications. The present section uses that bridge to state the
confirmation criterion and the stabilization consequence for physically non-selective theories.

\begin{theorem}[Reference-stabilization necessity]
\label{thm:reference-stabilization-necessity}
\label{cor:stabilization-requirement}
If a physically non-selective theory is used to underwrite theory-level confirmation and truth-apt
diachronic discourse whose measurement records function as evidential anchors, then the theory or
its admissible interpretive completion must restrict admissible reference assignments strongly enough
that fixation commutes with admissible redescription on the target discourse.
\end{theorem}

\begin{proof}
By Lemma~\ref{lem:dichotomy-admissible-assignments}, admissible assignments either collapse into a
single standing-bearing class on the target discourse or they do not. If they do not, then by
Theorem~\ref{thm:theory-level-confirmation-criterion} and
Theorem~\ref{thm:fixation-commutation-necessity}, theory-level confirmation through measurement
records is undefined on the fixed scope. But the standing conditional of
Section~\ref{sec:targetclass} assumes that the theory is being used to underwrite theory-level
confirmation and truth-apt diachronic discourse in which measurement records function as evidential
anchors. Therefore the role-divergent case is incompatible with the theory's intended explanatory use.
The only remaining option is that the theory or its admissible interpretive completion restricts
admissible reference assignments strongly enough for fixation to commute with admissible
redescription on the target discourse.
\end{proof}

\begin{proposition}[Selector and semantic-infrastructure normal form]
\label{prop:selector-normal-form}
Any successful stabilization strategy for admissible measurement-record reference takes exactly one of
two normal forms:
\begin{enumerate}
    \item \emph{selector form}, in which singular reference is fixed by selecting one history or one
    successor-relation as privileged; or
    \item \emph{semantic-infrastructure form}, in which admissible assignments remain multiple but are
    restricted strongly enough to collapse into one standing-equivalence class on the target
    discourse.
\end{enumerate}
No third same-scope stabilization form remains.
\end{proposition}

\begin{proof}
If a successful stabilization strategy fixes a unique referent outright, it has selector form. If it does
not fix a unique referent outright but still preserves theory-level confirmation, then by
Theorem~\ref{thm:reference-stabilization-necessity} it must restrict admissible assignments strongly
enough that fixation commutes and the assignments collapse into one standing-equivalence class on the
target discourse. That is semantic-infrastructure form. By
Proposition~\ref{prop:substrate-consequences}, no third same-scope locus of admissibility-bearing
work remains.
\end{proof}

\begin{proposition}[Scope of the main theorem]
\label{prop:primary-scope-confirmation}
Theorems~\ref{thm:theory-level-confirmation-criterion} and
\ref{thm:reference-stabilization-necessity} concern theory-level confirmation on the fixed scope. They
do not deny that branch-internal discourse may be perfectly determinate once a branch is taken as
fixed.
\end{proposition}

\begin{proof}
The hypotheses of Definitions~\ref{def:theory-level-confirmation} and
\ref{def:fixation-commutation} explicitly concern theory-level evidential use on the fixed downstream
substrate. Nothing in those definitions denies local branch-internal determinacy once a particular
branch-history is held fixed. The theorems only state that such local determinacy does not by itself
suffice for theory-level confirmation on the fixed scope.
\end{proof}
